\section{Introduction}\label{sec:introduction}

These notes assume a course in intermediate macroeconomics. We take national
income accounting, basic growth theory, and aggregate demand and supply as
familiar. The next step is to explain how aggregate outcomes follow from
decisions made over time, and how those decisions change with the economic
environment and policy.

The first three parts cover growth, fluctuations, and the open economy.
Part IV studies money, finance, and policy; Part V applies these mechanisms to
the dollar system and China's growth after the pandemic. Chapters marked
\textit{planned} contain outlines for later development.

\subsection{Growth, fluctuations, and their interaction}

Growth theory asks how productive capacity and living standards evolve. Capital
accumulation is the starting point. Saving decisions, population, technology,
and the allocation of resources then explain why accumulation and income differ
across economies.

Business-cycle theory asks how an economy responds to disturbances. A change in
productivity, demand, or financing conditions affects consumption, investment,
employment, and prices. The task is to identify the shock and explain the
adjustment that follows. A mechanism that amplifies a shock need not be the
mechanism that caused it.

The questions overlap. A recession can reduce investment and change future
capital. Population and productivity growth can change the conditions for
stabilization policy. Growth need not be an unaffected trend, and fluctuations
can have lasting effects. We begin with separate models to isolate mechanisms,
then study their interaction.

\subsection{From aggregate relations to individual decisions}

An aggregate consumption function such as $C=C(Y-T)$ summarizes a relationship.
It does not yet explain how a household responds to an expected tax change,
uncertain income, or a borrowing limit. To study those questions, we specify
what the household chooses and what constrains it.

\textbf{A two-period example.}
Let a household receive real income $y_0$ today and $y_1$ tomorrow. It starts
with no assets and leaves none after tomorrow. Saving $a_1$ earns a known gross
real return $R>0$. Negative $a_1$ is borrowing at the same return. The problem is
\begin{equation}\label{eq:intro-household}
\begin{split}
\max_{c_0,c_1,a_1}\;&u(c_0)+\beta u(c_1),\qquad 0<\beta<1,\\
\text{s.t. }&c_0+a_1=y_0,\\
&c_1=y_1+Ra_1.
\end{split}
\end{equation}
Assume $u'>0$, $u''<0$, positive consumption, and an interior solution. Substituting
the constraints leaves one choice:
\begin{equation}
\max_{a_1}\;u(y_0-a_1)+\beta u(y_1+Ra_1).
\end{equation}
Differentiate with respect to $a_1$, then substitute back for consumption:
\begin{align}
-u'(y_0-a_1)+\beta R u'(y_1+Ra_1)&=0,\\
u'(c_0)&=\beta R u'(c_1).\label{eq:intro-euler-certain}
\end{align}
Saving one more unit costs $u'(c_0)$ units of utility today and yields
$\beta R u'(c_1)$ tomorrow. An interior optimum equates the two.

For marginal utility $u'(c)=c^{-\sigma}$, $\sigma>0$,
\begin{align}
c_0^{-\sigma}&=\beta R c_1^{-\sigma},\\
\left(\frac{c_1}{c_0}\right)^\sigma&=\beta R,\\
\frac{c_1}{c_0}&=(\beta R)^{1/\sigma}.
\end{align}
Thus consumption growth depends on patience and the return to saving. Its
sensitivity to that return is governed by $1/\sigma$, the elasticity of
intertemporal substitution in this example. A binding borrowing constraint
would require an inequality condition instead of this interior equality.

\textbf{Expectations and policy.}
With uncertainty, the corresponding interior condition for an asset with gross
real payoff $R_{t+1}$ is
\begin{equation}\label{eq:intro-euler-uncertain}
u'(c_t)=\beta E_t\!\left[R_{t+1}u'(c_{t+1})\right],
\qquad E_t[\cdot]\equiv E[\cdot\mid\mathcal I_t].
\end{equation}
Here $\mathcal I_t$ is information available when the household chooses saving.
An uncertain return stays inside the expectation; its comovement with marginal
utility matters. Rational expectations require forecasts consistent with the
model and available information, not knowledge of tomorrow's realization.

The same logic matters for policy. An announced future tax change may alter
consumption today. A change in a policy rule can change expectations and the
aggregate relationships observed under the old rule. This is the concern in
the Lucas critique.\footnote{Robert E. Lucas, Jr. (1976), ``Econometric Policy
Evaluation: A Critique,'' \textit{Carnegie-Rochester Conference Series on Public
Policy}, 1, 19--46. \url{https://doi.org/10.1016/S0167-2231(76)80003-6}.}
Writing down individual decisions makes these responses explicit. It does not,
by itself, guarantee that the chosen preferences, constraints, or information
assumptions remain appropriate under every policy change.

\subsection{Reading a dynamic equilibrium}

The household takes income and returns as given. In general equilibrium, they
must be consistent with production and other agents' choices. A competitive
equilibrium combines individual optimality, budget constraints, market clearing,
and policy rules at each date and, under uncertainty, each history.

\textbf{States, choices, and paths.}
Consider a deterministic economy with fixed labor normalized to one, output
$f(k_t)$, and no government or foreign trade. Capital $k_t$ is available at the
start of period $t$. With depreciation $\delta$, feasibility gives
\begin{equation}\label{eq:intro-capital}
k_{t+1}=f(k_t)+(1-\delta)k_t-c_t.
\end{equation}
Capital is inherited; consumption is chosen and affects tomorrow's capital.
Here $k_t$ is a state variable and $c_t$ a control variable. The set of states
depends on the model's timing and assumptions.

If households can save in capital, competitive factor pricing gives the gross
return $f'(k_{t+1})+1-\delta$. Substituting this return into
\eqref{eq:intro-euler-certain} yields the intertemporal condition
\begin{equation}\label{eq:intro-growth-euler}
u'(c_t)=\beta\left[f'(k_{t+1})+1-\delta\right]u'(c_{t+1}).
\end{equation}
Together, \eqref{eq:intro-capital} and \eqref{eq:intro-growth-euler} describe
feasible accumulation and optimal consumption growth. Given $k_0$, different
initial choices of $c_0$ generate different candidate paths. Feasibility and
the Euler equation alone do not establish that each path is an equilibrium.
The model also needs its boundary conditions, including the appropriate
transversality condition for this infinite-horizon saving problem. We derive
these conditions in the Ramsey chapter.

\textbf{A steady state is a particular path.}
For a positive stationary allocation, set $k_{t+1}=k_t=k^*$ and
$c_{t+1}=c_t=c^*$ in the two equations. Since $u'(c^*)>0$,
\begin{align}
c^*&=f(k^*)-\delta k^*,\\
1&=\beta\left[f'(k^*)+1-\delta\right].
\end{align}
These equations locate a steady state. They do not describe the transition from
an arbitrary $k_0$, or prove existence or uniqueness of an equilibrium path.
With sustained technological or population growth, aggregate quantities may
keep growing on a balanced growth path. A stationary system may then require
normalization, for example by effective labor, rather than constant levels.

\textbf{Equilibrium and welfare answer different questions.}
A planner chooses a feasible allocation directly; a competitive equilibrium
also specifies prices that support individual choices. Under suitable
conditions, a planner's allocation can be decentralized. Taxes, market power,
externalities, incomplete markets, or nominal rigidities may prevent that
equivalence. Market clearing alone establishes neither efficiency nor the
desirability of a policy. A welfare comparison must state whose utility is
evaluated, which constraints apply, and what alternative is feasible.

\subsection{Accounting and notation used in these notes}

We need only a few accounting conventions at the outset. In a closed economy,
using gross investment $I_t$,
\begin{align}
Y_t&=C_t+I_t+G_t,\label{eq:intro-resource}\\
K_{t+1}&=(1-\delta)K_t+I_t.\label{eq:intro-stock-flow}
\end{align}
$K_t$ is the stock available at the beginning of the period; $I_t$ is a flow
during the period. Depreciation already enters \eqref{eq:intro-stock-flow}.
We do not add it again to \eqref{eq:intro-resource}. A household's purchase of
a financial claim need not create physical capital; domestic claims and
liabilities must be consolidated when deriving aggregate feasibility.

Accounting becomes a model only after we specify behavior, technology, and
equilibrium. Government budgets enter with fiscal policy; international
accounts enter in Part III.

Capital letters generally denote aggregates. Lower-case letters denote
per-person or normalized quantities where the chapter defines them. Some
existing chapters use local notation; the normalization must be checked before
comparing their equations. Unless stated otherwise, $\beta$ is a discount
factor, $\rho$ a rate of time preference, $r$ a net real interest rate, and
$R=1+r$ a gross real return. In discrete time, the convention
$\beta=1/(1+\rho)$ differs from continuous-time discounting $e^{-\rho t}$.

For a nominal one-period return $1+i_t$, realized inflation
$\pi_{t+1}=P_{t+1}/P_t-1$ gives the gross real return
\begin{equation}
R_{t+1}=\frac{1+i_t}{1+\pi_{t+1}}.
\end{equation}
An interest rate quoted in nominal units is therefore not the real return in
\eqref{eq:intro-euler-uncertain}. Under uncertainty, the realized price change
must remain inside the relevant expectation.

Finally, for a positive variable $X_t$ and a positive stationary reference
$\bar X$, define
\begin{equation}
\hat x_t\equiv\log X_t-\log\bar X
\approx\frac{X_t-\bar X}{\bar X}.
\end{equation}
The definition is exact; the approximation is local. A deviation from a steady
state, a growth rate, and an output gap relative to a flexible-price allocation
are different objects. Each chapter states its reference point before using
an approximation.
